\begin{figure}[t]
\centering
\fbox{
\begin{minipage}{0.95\linewidth}
\textbf{prompt}: there used to be a program for winmx called mcma, what does mcma stand for and what does the program do?
\vspace{1em}

\textbf{chosen}: MCMA stood for "Metis Chatroom Management Assistant." It was an add-on program for WinMX, which was a peer-to-peer file-sharing software popular in the early 2000s. WinMX allowed users to connect to chat rooms and share files with each other. \\
MCMA provided chat room administrators with tools to manage their chat rooms on the WinMX network. It included features like user kick/ban capabilities, word filtering, and automated messages. Essentially, it helped moderators to maintain order in the chat rooms and enforce the rules. 
WinMX and its associated tools, like MCMA, became less prevalent after the WinMX service was shut down due to legal pressures in 2005. Since then, many users have migrated to other file-sharing and chat platforms. Remember that using such tools for copyright infringement is illegal and against the terms of service of most platforms. Always ensure that you are following the law and the terms of service of any software or service you use.

\vspace{1em}

\textbf{rejected}: MCMA stands for Multicast Control Messages Analyzer and it was a program used to analyze winmx traffic. It was used to detect and log winmx traffic, allowing winmx users to see how much traffic they were sending and receiving. It also allowed users to monitor their connections and identify any potential problems.
\end{minipage}
}
\caption{Example 1 from the AlpacaEval Easy subset, which was rejected due to major hallucinations in both chosen and rejected responses.}
\label{fig:alpacaeval_1}
\end{figure}
